\originalsection{弹簧链、连续极限与场变分}
连续场不是把质点坐标随意换成函数. 原课从一条离散弹性链建立这种联系. 相邻质点平衡间隔为$a$, 第$i$个点的纵向位移为$\eta_i(t)$. 设线密度为$\lambda$, 每点质量$m=\lambda a$, 弹簧劲度$k=C/a$, 其中$C$为轴向刚度（弹性杆中可取Young模量乘截面积）. 离散拉格朗日量和方程为
\[
L_a=\sum_i\left[\frac{\lambda a}{2}\dot\eta_i^2
-\frac{C}{2a}(\eta_{i+1}-\eta_i)^2\right],\qquad
\lambda\ddot\eta_i=C\frac{\eta_{i+1}-2\eta_i+\eta_{i-1}}{a^2}.
\]
当位移变化的波长远大于$a$, 写$\eta_i(t)=\eta(x_i,t)$, 差商趋于$\eta_x$, 二阶差分趋于$\eta_{xx}$, $a\sum_i$趋于$\int\dd x$. 保持$\lambda,C$固定, 得
\[
L=\int\mathscr L\dd x,\qquad
\mathscr L=\frac\lambda2\eta_t^2-\frac C2\eta_x^2,
\qquad \lambda\eta_{tt}-C\eta_{xx}=0.
\]
波速$c=\sqrt{C/\lambda}$. 这里连续极限描述长波近似, 不保证有限$a$的链对任意短波也有相同色散. 例如离散平面波$\eta_i\propto e^{\mathrm i(kia-\omega t)}$给$\omega^2=4C\sin^2(ka/2)/(\lambda a^2)$, 仅在$|ka|\ll1$时接近$c^2k^2$.

同一场方程也可由驻定作用量得到. 对$S=\int\!\int\mathscr L(\eta,\eta_t,\eta_x,x,t)\dd x\dd t$, 在时间端点和固定空间边界取$\delta\eta=0$, 分别对$t,x$分部积分:
\[
\delta S=\int\!\int
\left(\mathscr L_\eta-\partial_t\mathscr L_{\eta_t}
-\partial_x\mathscr L_{\eta_x}\right)\delta\eta\dd x\dd t.
\]
因此场Euler--Lagrange方程为$\partial_t\mathscr L_{\eta_t}+\partial_x\mathscr L_{\eta_x}-\mathscr L_\eta=0$. 弹性链的密度给$\mathscr L_{\eta_t}=\lambda\eta_t$, $\mathscr L_{\eta_x}=-C\eta_x$, 立即重得波动方程. 原来每点一个方程, 现在每个空间标签$x$仍有一个方程, 空间导数表示邻点耦合. 三维时把$\partial_x\mathscr L_{\eta_x}$换成各空间方向之和; 多分量场则对每个分量独立变分. 这就是从有限自由度力学进入经典场论的具体入口.
